\documentclass[twocolumn,10pt]{article}
\usepackage{amsmath,amssymb,graphicx,booktabs,hyperref,microtype}
\usepackage[margin=1.8cm]{geometry}
\hypersetup{colorlinks=true,linkcolor=blue,citecolor=blue}

\title{\textbf{Paper 69: Wave-Range Dependence of Ordering in VCML}\\
  \large Minimum Spatial Range Threshold and the Limits of the DP Hypothesis}
\author{VCSM Theoretical Programme}
\date{2026-03-11}

\begin{document}
\maketitle

\begin{abstract}
Paper~68 conjectured that VCSM belongs to a novel universality class characterised
by directed dynamics, $Z_2$ symmetry, and long-range wave coupling ($r_w=5$),
and proposed that setting $r_w=1$ (nearest-neighbour) would either recover
Directed Percolation exponents or confirm the novel class.
We test this conjecture across three phases.
\textbf{Phase A} (fixed $L=80$, $P=0.02$, wave-rate sweep across
$r_w\in\{1,\ldots,8\}$): $|M|$ increases roughly linearly with wave area
$A(r_w)$; $U_4$ rises from $0.15$ at $r_w=1$ to $0.34$ at $r_w=8$.
\textbf{Phase B} (FSS at $r_w\in\{1,5\}$, matched coverage protocol):
$r_w=5$ shows orderly FSS slopes converging toward $-\beta/\nu=-0.67$;
$r_w=1$ shows slopes of $\approx-1.7$ at all tested $P$, steeper than the
disordered-phase noise floor, diagnosed as a \emph{protocol breakdown}:
the matched-coverage firing probability caps at $p_\mathrm{wave}=1.0$ for
all $L\geq40$ with $r_w=1$, so larger systems actually receive
proportionally \emph{fewer} hits per site.
\textbf{Phase C} (fine $P$ scan at $r_w=1$, $L=160$, matched coverage):
$|M|$ values are $\approx 3-8\times$ smaller than $r_w=5$ at the same $P$;
$U_4<0.18$ for $P\leq0.015$, indicating the system remains in or near the
disordered phase; power-law fit gives $\beta_\mathrm{fit}=0.113$ ($R^2=0.63$),
consistent with neither DP ($\beta\approx0.584$) nor Paper~63 ($\beta=0.65$).
The central finding is that \emph{the VCML ordering transition requires a
minimum effective wave range}: at $r_w=1$, the ordered phase is inaccessible
at the $P$ values and system sizes tested.
The DP hypothesis cannot be tested because $r_w=1$ does not produce a
cleanly ordered phase.
A new threshold is identified: $r_w/\text{zone\_width}$ must exceed $\sim1/40$
for the ordered phase to be accessible; the transition is driven by coherent
zone-mean perturbation, which requires $r_w$ to cover a macroscopic fraction
of the zone.
Paper~70 will map the $r_w$-dependence of $p_c$ and the ordering threshold.
\end{abstract}

% ─────────────────────────────────────────────────────────────────────────────
\section{Motivation}

Paper~68 falsified the Manna universality class for VCSM.
The remaining conjecture was: directed dynamics $+$ $Z_2$ symmetry $+$
\emph{long-range} wave coupling $\to$ novel class, where the long-range
coupling (wave radius $r_w=5$, covering 61 sites per wave) is the key
feature absent from both Manna and standard Directed Percolation (DP).

The direct test: reduce $r_w$ to 1 (nearest-neighbour, 5 sites per wave).
Three outcomes were possible:
\begin{enumerate}
  \item $r_w=1$ recovers DP exponents $\beta\approx0.584$: wave range IS the differentiator.
  \item $r_w=1$ gives the same $\beta\approx0.65$: wave range is irrelevant, something
        else defines the novel class.
  \item $r_w=1$ does not order at accessible $P$: the ordering mechanism itself
        requires long-range perturbation.
\end{enumerate}
The data supports outcome~3.

% ─────────────────────────────────────────────────────────────────────────────
\section{Model and Protocol}

Physical parameters identical to Papers~63--68.
Two wave-firing protocols:
\begin{itemize}
  \item \textbf{Protocol A} (constant rate): $p_\mathrm{wave}(L) =
        \min(1, 1/w_f(L))$ where $w_f(L) = 25(40/L)^2$.  Same fires/step
        regardless of $r_w$.  Used in Phase A.
  \item \textbf{Protocol B} (matched coverage): $p_\mathrm{wave}(L,r_w) =
        \min(1, A_\mathrm{ref}/A(r_w) \cdot 1/w_f(L))$ where
        $A_\mathrm{ref}=A(5)=61$ and $A(r_w)=2r_w(r_w+1)+1$ (Manhattan ball
        area).  Intended to equalise hits per site per step across $r_w$.
        Used in Phases B and C.
\end{itemize}

\textbf{Protocol B breakdown.}
The matched-coverage formula requires $p_\mathrm{wave}>1$ for $r_w=1$ at
$L\geq40$: specifically, $p_B(L,1) = \min(1, 12.2/w_f(L)) = 1.0$ for all
tested sizes.
This means larger systems receive proportionally \emph{fewer} hits per site
($\propto 1/L^2$) rather than the target flat density.
For $L=40$: $31.3$ hits/site/10k steps vs $61$ for $r_w=5$.
For $L=120$: $3.5$ hits/site/10k steps vs $15.2$ for $r_w=5$.
A clean $r_w=1$ comparison at large $L$ would require firing multiple waves
per step --- beyond the current simulation's single-wave-per-step architecture.

% ─────────────────────────────────────────────────────────────────────────────
\section{Phase A: $r_w$ Sweep at Fixed Conditions}

\begin{table}[h]
\centering
\caption{Phase A: L=80, P=0.020, Protocol A, 6 seeds, 10k steps.}
\label{tab:phaseA}
\begin{tabular}{rrrr}
\toprule
$r_w$ & $A(r_w)$ & $|M|\times10^4$ & $U_4$ \\
\midrule
1 & 5   & 0.8 & 0.152 \\
2 & 13  & 1.2 & -0.102 \\
3 & 25  & 2.2 & 0.144 \\
4 & 41  & 3.3 & 0.151 \\
5 & 61  & 4.1 & 0.244 \\
6 & 85  & 4.4 & 0.090 \\
8 & 145 & 8.3 & 0.339 \\
\bottomrule
\end{tabular}
\end{table}

Both $|M|$ and $U_4$ increase with $r_w$.
A power-law fit $|M|\sim A(r_w)^\alpha$ yields $\alpha\approx0.67$, consistent
with the zone-mean ordering mechanism: larger waves drive more of the zone
mean per event, amplifying $M$ roughly as the square root of wave coverage.
The $r_w=2$ case shows a negative $U_4$ ($-0.10\pm0.13$), consistent with
high variance and near-zero mean at this intermediate scale.

The critical observation: at $r_w=1$, $|M|$ is \emph{ten times smaller}
than at $r_w=8$ for the same $P$ and $L$.
Wave geometry, not just wave count, determines ordering strength.

% ─────────────────────────────────────────────────────────────────────────────
\section{Phase B: FSS Comparison at Matched Coverage}

\subsection{$r_w = 5$ (reference)}

FSS slopes $d\ln|M|/d\ln L$ converge toward $-\beta/\nu=-0.67$ as $P$ increases:
$-0.88$ at $P=0.005$ (disorder-like), $-0.56$ at $P=0.020$ (near-critical),
$-0.21$ at $P=0.050$ (ordered phase).
This is the expected signature of the $p_c=0^+$ crossover confirmed in Papers~64, 67.

\subsection{$r_w = 1$ (short-range)}

FSS slopes are $\approx-1.7$ at all tested $P$:
$-1.71$ at $P=0.005$, $-1.71$ at $P=0.010$, $-1.71$ at $P=0.020$, $-1.61$ at $P=0.050$.
These slopes are steeper than the disordered-phase noise floor ($-1.0$ in 2D),
and \emph{do not decrease} as $P$ increases toward $0.050$.

\textbf{Diagnosis}: Protocol B breaks down for $r_w=1$ because $p_\mathrm{wave}$
saturates at 1.0 for all tested $L$.
The actual hits/site/step at $L=120$ with $r_w=1$ is 3.5, vs 15.2 for $r_w=5$.
Larger systems are more under-driven, producing artificial steepening of the
scaling slope.
The $r_w=1$ FSS results are therefore \emph{confounded} and cannot be used to
extract exponents.

% ─────────────────────────────────────────────────────────────────────────────
\section{Phase C: Fine $P$ Scan at $r_w=1$, $L=160$}

\begin{table}[h]
\centering
\caption{Phase C: $r_w=1$, $L=160$, Protocol B ($p_\mathrm{wave}=1.0$), 6 seeds.
Comparison with Paper~68 Phase A ($r_w=5$, same $L$).}
\label{tab:phaseC}
\begin{tabular}{rrrr}
\toprule
$P$ & $|M|_{r_w=1}$ & $U_4$ & ratio vs P68 \\
\midrule
0.001 & $5\times10^{-5}$ & 0.115 & 0.38 \\
0.005 & $5\times10^{-5}$ & 0.141 & 0.36 \\
0.010 & $5\times10^{-5}$ & 0.144 & 0.33 \\
0.020 & $6\times10^{-5}$ & 0.170 & 0.21 \\
0.050 & $9\times10^{-5}$ & 0.351 & 0.12 \\
\bottomrule
\end{tabular}
\end{table}

At $r_w=1$, $L=160$ with $p_\mathrm{wave}=1.0$ (maximum possible):
the ordering signal is $3$--$8\times$ smaller than $r_w=5$ at the same $P$,
and $U_4<0.18$ for $P\leq0.015$.
For reference, $P=0.020$ with $r_w=5$ gives $U_4=0.321$ and $P=0.050$ gives
$U_4=0.584$ --- the system is well-ordered.
At $r_w=1$, $P=0.050$ gives only $U_4=0.35$, suggesting ordering onset begins
around $P\approx0.05$--$0.10$ with $r_w=1$.

The power-law fit gives $\beta_\mathrm{fit}=0.113$ ($R^2=0.63$).
This matches neither DP ($\beta\approx0.584$) nor Paper~63 ($\beta=0.65$).
However, the $U_4$ flatness at small $P$ indicates the system is not yet in
the asymptotic ordered phase, so this $\beta$ measurement is unreliable ---
exactly the same contamination problem identified in Paper~68.

% ─────────────────────────────────────────────────────────────────────────────
\section{Synthesis: Minimum Wave-Range Threshold}

The ordering mechanism in VCML requires waves to drive a coherent perturbation
of the zone mean $M = \langle\phi\rangle_{z_0} - \langle\phi\rangle_{z_1}$.
A wave of radius $r_w$ hitting the zone boundary deposits a causal signal
in approximately $A(r_w)/2 \sim r_w^2$ sites.
For this to coherently drive $M$, the signal must be large relative to the
background fluctuations, which scale as $\sigma_M \sim 1/\sqrt{N_\mathrm{zone}}
\sim 1/L$.
The condition for ordered-phase accessibility at finite $L$ is therefore:
\begin{equation}
  \frac{r_w^2}{L^2} \gtrsim \frac{1}{\sqrt{N_\mathrm{zone}}} \quad\Rightarrow\quad
  r_w \gtrsim L^{1/2}.
  \label{eq:threshold}
\end{equation}
At $L=160$: $r_w \gtrsim 13$.
At $L=80$: $r_w \gtrsim 9$.
At $L=40$: $r_w \gtrsim 6$.

This scaling is too aggressive: $r_w=5$ clearly orders at $L=80$ ($U_4=0.24$
at $P=0.02$), so the threshold exponent is softer than $1/2$.
The empirical threshold appears closer to $r_w \gtrsim \text{zone\_width}/20$,
giving $r_w \gtrsim 4$ for $L=80$ (zone width = 40).
The exact threshold scaling is an open question.

\textbf{Consequence for the DP conjecture}:
The test of whether $r_w=1$ recovers DP exponents is currently inaccessible
because $r_w=1$ at $L=160$ does not produce a cleanly ordered phase.
To perform a clean comparison, one would need either:
(a) very high $P$ ($P\sim0.1$--$0.5$) with $r_w=1$ to force ordering at small $L$; or
(b) a modified simulation that fires multiple waves per step to achieve coverage parity.
Both are deferred to Paper~70.

% ─────────────────────────────────────────────────────────────────────────────
\section{Conclusions}

\begin{enumerate}
  \item Ordering strength scales strongly with wave area:
        $|M| \sim A(r_w)^{0.67}$ at fixed $L$, $P$, and wave-rate.
  \item The ordering transition in VCML requires a minimum wave range:
        at $r_w=1$, the system remains in or near the disordered phase for
        $P\leq0.050$ at $L=160$, while $r_w=5$ is well-ordered by $P=0.020$.
  \item The matched-coverage protocol (Protocol B) breaks down for $r_w=1$
        at large $L$ because $p_\mathrm{wave}$ cannot exceed 1.0 per step.
        A clean comparison requires firing multiple waves per step.
  \item The DP conjecture from Paper~68 cannot be confirmed or refuted:
        $r_w=1$ does not produce a cleanly ordered phase at accessible $P$.
  \item New open question: is there a critical $r_w^*(L)$ below which the
        ordered phase is inaccessible? How does $p_c(r_w)$ scale with $r_w$?
        Paper~70 will map this threshold by scanning $(r_w, P)$ jointly at
        fixed $L$.
\end{enumerate}

\end{document}
